\begin{proposition}\label{prop:nonic-no62}
There is no point of order $62$ on an elliptic curve over a number
field of degree $9$.
\end{proposition}
\begin{proof}
Put $Y=X_1(31)$ and $J=J_1(31)$. The curve has genus $26$ and
$\operatorname{gon}_{\mathbb Q}(Y)=12$ \cite{DvH}. By
\cite[Theorem~3.1(2), Theorem~4.13 and Table~2]{DEHMZ},
\[
 J(\mathbb Q)\simeq C_{10}\times C_{1772833370}.
\]
We use the degree-$3$ forgetful map $\pi:X_1(62)\to Y$.
Its noncuspidal fibre in characteristic $3$ is
$E[2]\setminus\{0\}$.

Suppose a point of order $62$ is defined over a nonic field.
Let $D$ be its trace divisor and put $D_Y=\pi_*D$.
The global and Frobenius-$3$ cusp orbit partitions on $Y$ both
have signature $1^{15}15$. Indeed, $3$ has order $15$ modulo
sign in $(\mathbb Z/31\mathbb Z)^\times$.
If $\overline D_Y$ were entirely cuspidal, it would lift to a
rational cuspidal divisor $C$ of degree $9$.
Since $J(\mathbb Q)$ is finite and torsion reduction is injective
at the good prime $3$, we would have $D_Y\sim C$.
These distinct rational effective divisors contradict the gonality
of $Y$. Thus the source reduction has a nonempty noncuspidal part.

The Hasse bound excludes $62$-torsion over $\mathbb F_{3^e}$ for
$e\le4$: for $e=4$ its interval is $[64,100]$.
The noncuspidal source part is therefore one closed point $Q_e$,
with multiplicity $1$ and $5\le e\le9$.
Let $P_d=\pi(Q_e)$ and $e=fd$.
For $f=2$ the only possibilities are $(d,e)=(3,6),(4,8)$.
Over $\mathbb F_{27}$ or $\mathbb F_{81}$ a curve with $31$-torsion
would have respectively $31$ or $93$ points, both odd.
Its $2$-division polynomial is an irreducible cubic, so a quadratic
fibre point is impossible. For $f=3$, one has $d=2$ or $3$.
The Hasse bound excludes $d=2$, and the exact point count on $Y$
gives no noncuspidal place of degree $3$.
Consequently $f=1$: we need only consider a place $P_d$ of degree
$5\le d\le9$ whose elliptic curve has a rational nonzero
$2$-torsion point.

The proven finite class group is
\[
 G=J(\mathbb F_3)\simeq C_2^2\times C_{10}\times C_{3352427902670},
 \qquad\#G=134097116106800.
\]
The rational image $H$ has index $7564=2^2\cdot31\cdot61$.
The class presentation uses all $3481$ places of degrees at most
$9$, attaining the unconditional generation bound $9$, and has
$4391$ principal relation rows.
The noncuspidal closed-place counts on $Y$ are
\[
 (0,0,0,15,91,120,330,930,1980),
\]
and those on $X_1(62)$ are
\[
 (0,0,0,0,180,150,330,930,1600).
\]
Independent finite-field arithmetic verifies every Tate normal
form, its exact $31$-torsion point, all $13$ generators of its
factor-basis ideal, distinct Frobenius orbits and both count lists.

Fix a rational cusp $c_0$.
Every cuspidal complement here has degree at most $4$, hence is
supported on rational cusps. It is necessary that $[P_d-dc_0]\in H$.
There are respectively $90,90,180,630,1200$ eligible places in
degrees $5,6,7,8,9$, giving $2190$ tests.
The complete relation matrix gives character images of orders
$16,31,61$ in targets of exponents $16,31,61$.
The rational cusps give images of orders $4,1,1$.
Lemma~\ref{lem:nonic35-primary-characters} gives an exact test for $H$.
Independent integer arithmetic replays all relation identities,
image orders and tests. Exactly five places survive, all of degree $9$.

It remains to lift these five places. Let $a$ satisfy
\[
 a^9-2a^8+a^6-a^5+14a^4-28a^3+19a^2-2a-1=0,
\]
and put
\[
 x=a,\qquad
 119y=9a^8-28a^7+16a^6+46a^5-11a^4+125a^3-440a^2+386a+150.
\]
This is the nonic $31$-torsion example of
Proposition~\ref{prop:nonic-sporadic-realizations}.
Its field polynomial is irreducible modulo $3$.
The fifteen diamond translates of its degree-$9$ trace reduce to
exactly the five surviving places. This identification uses exact
Tate normalization and Frobenius orbit keys.

For any surviving reduction, choose the corresponding global
trace divisor $B$. Then $\overline D_Y=\overline B$.
The rational torsion class $[D_Y-B]$ reduces to zero, so
$D_Y\sim B$. Since $9<\operatorname{gon}_{\mathbb Q}(Y)$, we obtain
$D_Y=B$ as divisors.
The point on $Y$ therefore has the same nonic residue field and
underlying elliptic curve as this diamond translate.
The curve has full torsion $C_{31}$ over that field, as proved in
Proposition~\ref{prop:nonic-sporadic-realizations}; in particular,
its good degree-$1$ reduction of order $93$ at $83$ excludes
rational $2$-torsion.
It cannot lift to $X_1(62)$ over the same nonic field.
This proves the proposition.
\end{proof}
